\documentclass[11pt,a4paper]{article}
\usepackage{fontspec,amsmath,amssymb,amsthm,microtype}
\setmainfont{lmroman10-regular.otf}[BoldFont=lmroman10-bold.otf,ItalicFont=lmroman10-italic.otf,BoldItalicFont=lmroman10-bolditalic.otf]
\usepackage{xcolor}
\usepackage[most]{tcolorbox}
\usepackage{mathrsfs}
\providecommand{\lt}{<}
\providecommand{\gt}{>}
\setmonofont{DejaVuSansMono.ttf}[BoldFont=DejaVuSansMono-Bold.ttf,ItalicFont=DejaVuSansMono-Oblique.ttf,BoldItalicFont=DejaVuSansMono-BoldOblique.ttf,Scale=0.86]
\newfontfamily\leanfallback{STIXTwoMath-Regular.otf}[Scale=0.86]
\definecolor{leanrule}{gray}{0.62}
\definecolor{leanback}{gray}{0.962}
\definecolor{leantitle}{gray}{0.22}
\definecolor{leanmuted}{gray}{0.38}
\newtcolorbox{leanbox}[1]{enhanced,breakable,colback=leanback,colframe=leanrule,
  boxrule=0.4pt,arc=0.6mm,left=2.4mm,right=2.4mm,top=1.8mm,bottom=1.8mm,
  before skip=0.8em,after skip=1em,fonttitle=\small\bfseries,coltitle=leantitle,
  colbacktitle=leanback,titlerule=0.3pt,title={#1}}
\newlength{\leancharwidth}
\newenvironment{leancode}{\par\smallskip\begingroup\ttfamily\footnotesize\raggedright
  \settowidth{\leancharwidth}{M}\setlength{\parskip}{0pt}\setlength{\parindent}{0pt}
  \frenchspacing}{\par\endgroup\smallskip}
\newcommand{\leanline}[2]{\hangindent=\dimexpr#1\leancharwidth+4\leancharwidth\relax
  \hangafter=1\noindent\hspace*{#1\leancharwidth}#2\par}
\newcommand{\leansource}[1]{\par{\footnotesize\color{leanmuted}#1}\par}
\newcommand{\leanlabel}[1]{\par\smallskip\noindent{\small\bfseries #1}\par}
\usepackage[margin=25mm]{geometry}
\usepackage[unicode,hidelinks]{hyperref}
\hypersetup{pdfauthor={LangevinLean Project}}
\setlength{\emergencystretch}{3em}
\setlength{\parskip}{0.4em}
\setlength{\parindent}{0pt}
\allowdisplaybreaks[2]
\raggedbottom
\title{\texorpdfstring{Particle-entropy convergence from the lifted logarithmic Sobolev inequality}{Particle-entropy convergence from the lifted logarithmic Sobolev inequality}}
\author{LangevinLean Project}
\date{}
\begin{document}\maketitle
Draft mathematical exposition; these notes may contain mistakes.


The Gaussian interpolation theorem derives the required density, differentiation, and integrability properties from the entropy assumptions. For the particle model, substitute \(V_N\), step \(\gamma/N\), gradient constant \(NL\), and dimension \(D=d_\theta+Nd_x\). The lifted-target logarithmic Sobolev inequality, positive admissible step, and finite initial joint entropy then imply finite output entropy and the quantitative estimate. Division by particle count and parameter projection give the stated corollaries.



Earlier results used here: \cite[Section~3.29]{LangevinLean2026ULALSI}.

\section{\texorpdfstring{Finite joint entropy and its convergence estimate}{Finite joint entropy and its convergence estimate}}


Let \(N>0\), \(D=d_\theta+Nd_x\), and \(V_N(z)=\sum_iU(\theta,\sqrt N z_i)\). Assume differentiability of \(U\) with \(L\)-Lipschitz gradient, a logarithmic Sobolev inequality of constant \(\beta>0\) for \(\Pi_N\), a positive physical step satisfying \(\gamma\le\beta/(4NL^2)\), and finite initial joint entropy \(H_0\). Every particle law then has finite divergence and satisfies

\[H_k\le e^{-\beta\gamma k/N}H_0+\frac{8\gamma ND L^2}{\beta}.\]

The lifted potential is differentiable with \(NL\)-Lipschitz gradient, and its Gaussian Euler law at step \(h=\gamma/N\) is the particle law. The physical restriction is equivalent to \(h\le\beta/(4(NL)^2)\). Thus the Gaussian interpolation theorem applies with \(h=\gamma/N\), dimension \(D\), and gradient constant \(NL\). Its rate is \(e^{-\beta\gamma k/N}\) and its error is \(8\gamma ND L^2/\beta\). The admissible positive step also gives positive \(L\) before any division.



In detail, the squared gradient constant is \(\ell^2=N^2L^2\) and the step is \(h=\gamma/N\). These substitutions give restriction \(\gamma\le\beta/(4NL^2)\) and error term \(8\gamma NDL^2/\beta\). The Gaussian theorem supplies the interpolation regularity and the compatibility of constants, so neither is an additional premise here. The finite-divergence conclusion is retained alongside the real inequality, justifying each real-valued entropy in the formula.

\begin{leanbox}{Lean: \texttt{particle\_\allowbreak{}kl\_\allowbreak{}bound\_\allowbreak{}of\_\allowbreak{}finite\_\allowbreak{}kl}}
\leansource{\texttt{particle\_\allowbreak{}kl\_\allowbreak{}bound\_\allowbreak{}of\_\allowbreak{}finite\_\allowbreak{}kl} \quad IPLA/\allowbreak{}LSI/\allowbreak{}UnconditionalTheorem.lean, lines 22–52}
\begin{leancode}
\leanline{0}{theorem\ particle\_kl\_bound\_of\_finite\_kl}
\leanline{4}{\{N\ dθ\ dx\ :\ ℕ\}\ (hN\ :\ 0\ <\ N)\ (U\ :\ Joint\ dθ\ dx\ →\ ℝ)}
\leanline{4}{(β\ γ\ :\ ℝ)\ (L\ :\ ℝ≥0)\ (μ\ :\ Measure\ (State\ N\ dθ\ dx))\ [IsProbabilityMeasure\ μ]}
\leanline{4}{(hdiff\ :\ Differentiable\ ℝ\ U)\ (hLip\ :\ LipschitzWith\ L\ (gradient\ U))}
\leanline{4}{(hβ\ :\ 0\ <\ β)\ (hγ\ :\ 0\ <\ γ)}
\leanline{4}{(hsmall\ :\ γ\ ≤\ β\ /\ (4\ *\ (N\ :\ ℝ)\ *\ (L\ :\ ℝ)\ \textasciicircum{}\ 2))}
\leanline{4}{(hLSI\ :\ ULA.LSI.IsLogSobolev\ (@gibbsTarget\ N\ dθ\ dx\ U)\ β)}
\leanline{4}{(hKL\ :\ InformationTheory.klDiv\ μ\ (gibbsTarget\ U)\ ≠\ ⊤)\ (k\ :\ ℕ)\ :}
\leanline{4}{InformationTheory.klDiv\ (particleLaws\ U\ γ\ μ\ k)\ (gibbsTarget\ U)\ ≠\ ⊤\ ∧}
\leanline{6}{(InformationTheory.klDiv\ (particleLaws\ U\ γ\ μ\ k)\ (gibbsTarget\ U)).toReal\ ≤}
\leanline{8}{Real.exp\ (-β\ *\ γ\ *\ k\ /\ N)\ *}
\leanline{10}{(InformationTheory.klDiv\ μ\ (gibbsTarget\ U)).toReal\ +}
\leanline{12}{8\ *\ γ\ *\ N\ *\ (dθ\ +\ N\ *\ dx\ :\ ℕ)\ *\ (L\ :\ ℝ)\ \textasciicircum{}\ 2\ /\ β}
\end{leancode}
\leanlabel{How the Lean proof follows the argument}
\texttt{particle\_\allowbreak{}kl\_\allowbreak{}bound\_\allowbreak{}of\_\allowbreak{}finite\_\allowbreak{}kl} derives \texttt{hn} from positive particle count and proves \texttt{h\allowbreak{}L} by excluding zero smoothness: that case contradicts the admissible positive step. \texttt{hs} and \texttt{hscaled} express the lifted step restriction. \texttt{ULA.\allowbreak{}LSI.\allowbreak{}ula\_\allowbreak{}kl\_\allowbreak{}bound} supplies \texttt{hfinite} and \texttt{hb} for the actual summed potential. \texttt{potential\_\allowbreak{}differentiable} and \texttt{potential\_\allowbreak{}gradient\_\allowbreak{}lipschitz} supply its analytic hypotheses; the remaining simplification uses \texttt{finrank\_\allowbreak{}state}, \texttt{lifted\_\allowbreak{}entropy\_\allowbreak{}rate\_\allowbreak{}eq} and \texttt{lifted\_\allowbreak{}entropy\_\allowbreak{}bias\_\allowbreak{}eq}. No regular-interpolation data is an additional premise here.
\end{leanbox}


Earlier results used here: \cite[Section~3.1.3]{LangevinLean2026ULALSI}; \cite[Section~3.29.7]{LangevinLean2026ULALSI}; \cite[Section~3.1.1]{LangevinLean2026IPLA}; \cite[Section~3.1.3]{LangevinLean2026IPLA}; \cite[Section~3.1.4]{LangevinLean2026IPLA}; \cite[Section~3.3.2]{LangevinLean2026IPLA}.

\section{\texorpdfstring{The entropy bound divided by particle count}{The entropy bound divided by particle count}}


Divide the finite joint entropy inequality by positive \(N\). Division preserves order and gives

\[\frac{H_k}{N}\le e^{-\beta\gamma k/N}\frac{H_0}{N}+\frac{8\gamma D L^2}{\beta}.\]

Because \(N>0\), every term is divided by the same finite positive number. The exponential is unchanged, and its initial factor is the initial joint entropy per particle. This is a numerical normalization of the joint divergence, not the divergence of the parameter marginal.

\begin{leanbox}{Lean: \texttt{particle\_\allowbreak{}kl\_\allowbreak{}bound\_\allowbreak{}normalized\_\allowbreak{}of\_\allowbreak{}finite\_\allowbreak{}kl}}
\leansource{\texttt{particle\_\allowbreak{}kl\_\allowbreak{}bound\_\allowbreak{}normalized\_\allowbreak{}of\_\allowbreak{}finite\_\allowbreak{}kl} \quad IPLA/\allowbreak{}LSI/\allowbreak{}UnconditionalTheorem.lean, lines 56–72}
\begin{leancode}
\leanline{0}{theorem\ particle\_kl\_bound\_normalized\_of\_finite\_kl}
\leanline{4}{\{N\ dθ\ dx\ :\ ℕ\}\ (hN\ :\ 0\ <\ N)\ (U\ :\ Joint\ dθ\ dx\ →\ ℝ)}
\leanline{4}{(β\ γ\ :\ ℝ)\ (L\ :\ ℝ≥0)\ (μ\ :\ Measure\ (State\ N\ dθ\ dx))\ [IsProbabilityMeasure\ μ]}
\leanline{4}{(hdiff\ :\ Differentiable\ ℝ\ U)\ (hLip\ :\ LipschitzWith\ L\ (gradient\ U))}
\leanline{4}{(hβ\ :\ 0\ <\ β)\ (hγ\ :\ 0\ <\ γ)}
\leanline{4}{(hsmall\ :\ γ\ ≤\ β\ /\ (4\ *\ (N\ :\ ℝ)\ *\ (L\ :\ ℝ)\ \textasciicircum{}\ 2))}
\leanline{4}{(hLSI\ :\ ULA.LSI.IsLogSobolev\ (@gibbsTarget\ N\ dθ\ dx\ U)\ β)}
\leanline{4}{(hKL\ :\ InformationTheory.klDiv\ μ\ (gibbsTarget\ U)\ ≠\ ⊤)\ (k\ :\ ℕ)\ :}
\leanline{4}{InformationTheory.klDiv\ (particleLaws\ U\ γ\ μ\ k)\ (gibbsTarget\ U)\ ≠\ ⊤\ ∧}
\leanline{6}{(InformationTheory.klDiv\ (particleLaws\ U\ γ\ μ\ k)\ (gibbsTarget\ U)).toReal\ /\ N\ ≤}
\leanline{8}{Real.exp\ (-β\ *\ γ\ *\ k\ /\ N)\ *}
\leanline{10}{((InformationTheory.klDiv\ μ\ (gibbsTarget\ U)).toReal\ /\ N)\ +}
\leanline{12}{8\ *\ γ\ *\ (dθ\ +\ N\ *\ dx\ :\ ℕ)\ *\ (L\ :\ ℝ)\ \textasciicircum{}\ 2\ /\ β}
\end{leancode}
\leanlabel{How the Lean proof follows the argument}
\texttt{particle\_\allowbreak{}kl\_\allowbreak{}bound\_\allowbreak{}normalized\_\allowbreak{}of\_\allowbreak{}finite\_\allowbreak{}kl} obtains the finite joint divergence \texttt{hfinite} and full bound \texttt{hb} from \texttt{particle\_\allowbreak{}kl\_\allowbreak{}bound\_\allowbreak{}of\_\allowbreak{}finite\_\allowbreak{}kl}. It keeps the former as the first component of its conclusion and applies \texttt{normalize\_\allowbreak{}entropy\_\allowbreak{}bound} to the latter. Thus every term of the full inequality, including the initial entropy, is divided by the same positive particle count.



The helper \texttt{normalize\_\allowbreak{}entropy\_\allowbreak{}bound} divides both the initial contribution and the bias, canceling one particle factor in the latter. The same finite-KL witness is retained for the unnormalized divergence. This quantity is the entropy of the full state divided by particle number; parameter projection is handled separately.
\end{leanbox}


Earlier results used here: \cite[Section~3.1.3]{LangevinLean2026ULALSI}; \cite[Section~3.1.1]{LangevinLean2026IPLA}; \cite[Section~3.3.2]{LangevinLean2026IPLA}.

\section{\texorpdfstring{A linearly scaled lifted logarithmic Sobolev constant}{A linearly scaled lifted logarithmic Sobolev constant}}


Assume the lifted Gibbs target has logarithmic Sobolev constant \(\beta=N\alpha\) with \(\alpha>0\). Positivity of particle count gives \(\beta>0\), and the restriction \(\gamma\le\beta/(4NL^2)\) becomes \(\gamma\le\alpha/(4L^2)\). Apply the normalized joint theorem with this constant. Its rate is \(e^{-\alpha\gamma k}\). The dimension equality \(D/N=d_x+d_\theta/N\) gives error term \((8\gamma L^2/\alpha)(d_x+d_\theta/N)\), so

\[\frac{H_k}{N}\le e^{-\alpha\gamma k}\frac{H_0}{N}+\frac{8\gamma L^2}{\alpha}\left(d_x+\frac{d_\theta}{N}\right).\]

The functional-inequality premise concerns \(\Pi_N\) itself. The relation \(D/N=d_x+d_\theta/N\) separates the intensive latent and parameter contributions. Since the particles share a parameter, no product tensorization of independent pair targets is used to obtain the assumed constant.

\begin{leanbox}{Lean: \texttt{particle\_\allowbreak{}kl\_\allowbreak{}bound\_\allowbreak{}linear\_\allowbreak{}lsi\_\allowbreak{}of\_\allowbreak{}finite\_\allowbreak{}kl}}
\leansource{\texttt{particle\_\allowbreak{}kl\_\allowbreak{}bound\_\allowbreak{}linear\_\allowbreak{}lsi\_\allowbreak{}of\_\allowbreak{}finite\_\allowbreak{}kl} \quad IPLA/\allowbreak{}LSI/\allowbreak{}UnconditionalTheorem.lean, lines 76–96}
\begin{leancode}
\leanline{0}{theorem\ particle\_kl\_bound\_linear\_lsi\_of\_finite\_kl}
\leanline{4}{\{N\ dθ\ dx\ :\ ℕ\}\ (hN\ :\ 0\ <\ N)\ (U\ :\ Joint\ dθ\ dx\ →\ ℝ)}
\leanline{4}{(α\ γ\ :\ ℝ)\ (L\ :\ ℝ≥0)\ (μ\ :\ Measure\ (State\ N\ dθ\ dx))\ [IsProbabilityMeasure\ μ]}
\leanline{4}{(hdiff\ :\ Differentiable\ ℝ\ U)\ (hLip\ :\ LipschitzWith\ L\ (gradient\ U))}
\leanline{4}{(hα\ :\ 0\ <\ α)\ (hγ\ :\ 0\ <\ γ)}
\leanline{4}{(hsmall\ :\ γ\ ≤\ α\ /\ (4\ *\ (L\ :\ ℝ)\ \textasciicircum{}\ 2))}
\leanline{4}{(hLSI\ :\ ULA.LSI.IsLogSobolev\ (@gibbsTarget\ N\ dθ\ dx\ U)\ ((N\ :\ ℝ)\ *\ α))}
\leanline{4}{(hKL\ :\ InformationTheory.klDiv\ μ\ (gibbsTarget\ U)\ ≠\ ⊤)\ (k\ :\ ℕ)\ :}
\leanline{4}{InformationTheory.klDiv\ (particleLaws\ U\ γ\ μ\ k)\ (gibbsTarget\ U)\ ≠\ ⊤\ ∧}
\leanline{6}{(InformationTheory.klDiv\ (particleLaws\ U\ γ\ μ\ k)\ (gibbsTarget\ U)).toReal\ /\ N\ ≤}
\leanline{8}{Real.exp\ (-α\ *\ γ\ *\ k)\ *}
\leanline{10}{((InformationTheory.klDiv\ μ\ (gibbsTarget\ U)).toReal\ /\ N)\ +}
\leanline{12}{(8\ *\ γ\ *\ (L\ :\ ℝ)\ \textasciicircum{}\ 2\ /\ α)\ *\ ((dx\ :\ ℝ)\ +\ (dθ\ :\ ℝ)\ /\ N)}
\end{leancode}
\leanlabel{How the Lean proof follows the argument}
\texttt{particle\_\allowbreak{}kl\_\allowbreak{}bound\_\allowbreak{}linear\_\allowbreak{}lsi\_\allowbreak{}of\_\allowbreak{}finite\_\allowbreak{}kl} proves positive smoothness using \texttt{ULA.\allowbreak{}LSI.\allowbreak{}positive\_\allowbreak{}smoothness\_\allowbreak{}of\_\allowbreak{}step}, then rewrites the step assumption with \texttt{linear\_\allowbreak{}lsi\_\allowbreak{}step\_\allowbreak{}iff}. It applies the normalized theorem at the supplied linearly scaled lifted-target constant. The final rewrites \texttt{linear\_\allowbreak{}entropy\_\allowbreak{}rate\_\allowbreak{}eq} and \texttt{linear\_\allowbreak{}normalized\_\allowbreak{}entropy\_\allowbreak{}bias\_\allowbreak{}eq} give the stated exponent and effective dimension while retaining finite joint divergence.
\end{leanbox}


Earlier results used here: \cite[Section~3.1.3]{LangevinLean2026ULALSI}; \cite[Section~3.29.1]{LangevinLean2026ULALSI}; \cite[Section~3.1.1]{LangevinLean2026IPLA}; \cite[Section~3.3.2]{LangevinLean2026IPLA}.

\section{\texorpdfstring{The parameter-entropy consequence}{The parameter-entropy consequence}}


Project both the particle law and its Gibbs reference to their parameter coordinates. Data processing gives \(\operatorname{KL}(\mu_k^\theta\Vert\Pi_N^\theta)\le H_k\). The joint theorem supplies \(H_k<\infty\), so the projected divergence is finite as well. Substitution of the full numerical estimate yields

\[\operatorname{KL}(\mu_k^\theta\Vert\Pi_N^\theta)\le e^{-\beta\gamma k/N}H_0+\frac{8\gamma ND L^2}{\beta}.\]

The initial quantity is the full joint entropy \(H_0\). The parameter divergence is unnormalized, so its error retains factor \(N\). If a joint bound is supplied per particle, it must first be multiplied by \(N\). Data processing changes only the left-hand divergence and does not replace the initial joint entropy by a smaller projected quantity.

\begin{leanbox}{Lean: \texttt{parameter\_\allowbreak{}kl\_\allowbreak{}bound\_\allowbreak{}of\_\allowbreak{}finite\_\allowbreak{}kl}}
\leansource{\texttt{parameter\_\allowbreak{}kl\_\allowbreak{}bound\_\allowbreak{}of\_\allowbreak{}finite\_\allowbreak{}kl} \quad IPLA/\allowbreak{}LSI/\allowbreak{}UnconditionalTheorem.lean, lines 100–121}
\begin{leancode}
\leanline{0}{theorem\ parameter\_kl\_bound\_of\_finite\_kl}
\leanline{4}{\{N\ dθ\ dx\ :\ ℕ\}\ (hN\ :\ 0\ <\ N)\ (U\ :\ Joint\ dθ\ dx\ →\ ℝ)}
\leanline{4}{(β\ γ\ :\ ℝ)\ (L\ :\ ℝ≥0)\ (μ\ :\ Measure\ (State\ N\ dθ\ dx))\ [IsProbabilityMeasure\ μ]}
\leanline{4}{(hdiff\ :\ Differentiable\ ℝ\ U)\ (hLip\ :\ LipschitzWith\ L\ (gradient\ U))}
\leanline{4}{(hβ\ :\ 0\ <\ β)\ (hγ\ :\ 0\ <\ γ)}
\leanline{4}{(hsmall\ :\ γ\ ≤\ β\ /\ (4\ *\ (N\ :\ ℝ)\ *\ (L\ :\ ℝ)\ \textasciicircum{}\ 2))}
\leanline{4}{(hLSI\ :\ ULA.LSI.IsLogSobolev\ (@gibbsTarget\ N\ dθ\ dx\ U)\ β)}
\leanline{4}{(hKL\ :\ InformationTheory.klDiv\ μ\ (gibbsTarget\ U)\ ≠\ ⊤)\ (k\ :\ ℕ)\ :}
\leanline{4}{InformationTheory.klDiv\ ((particleLaws\ U\ γ\ μ\ k).map\ parameterCLM)}
\leanline{6}{(@parameterTarget\ N\ dθ\ dx\ U)\ ≠\ ⊤\ ∧}
\leanline{6}{(InformationTheory.klDiv\ ((particleLaws\ U\ γ\ μ\ k).map\ parameterCLM)}
\leanline{8}{(@parameterTarget\ N\ dθ\ dx\ U)).toReal\ ≤}
\leanline{8}{Real.exp\ (-β\ *\ γ\ *\ k\ /\ N)\ *}
\leanline{10}{(InformationTheory.klDiv\ μ\ (gibbsTarget\ U)).toReal\ +}
\leanline{12}{8\ *\ γ\ *\ N\ *\ (dθ\ +\ N\ *\ dx\ :\ ℕ)\ *\ (L\ :\ ℝ)\ \textasciicircum{}\ 2\ /\ β}
\end{leancode}
\leanlabel{How the Lean proof follows the argument}
\texttt{parameter\_\allowbreak{}kl\_\allowbreak{}bound\_\allowbreak{}of\_\allowbreak{}finite\_\allowbreak{}kl} installs probability instances for the numerical and target laws, then obtains \texttt{hfinite} and \texttt{hb} from the full joint entropy theorem. \texttt{parameter\_\allowbreak{}kl\_\allowbreak{}ne\_\allowbreak{}top\_\allowbreak{}of\_\allowbreak{}gibbs\_\allowbreak{}joint} supplies finite marginal divergence and \texttt{parameter\_\allowbreak{}kl\_\allowbreak{}bound\_\allowbreak{}of\_\allowbreak{}gibbs\_\allowbreak{}joint\_\allowbreak{}bound} transfers the full bound through the parameter projection. It is the unnormalized joint estimate that is passed to data processing.
\end{leanbox}


Earlier results used here: \cite[Section~3.1.3]{LangevinLean2026ULALSI}; \cite[Section~3.1.1]{LangevinLean2026IPLA}; \cite[Section~3.1.6]{LangevinLean2026IPLA}; \cite[Section~3.3.2]{LangevinLean2026IPLA}.

\begin{thebibliography}{ULA-LSI}
\bibitem[ULA-LSI]{LangevinLean2026ULALSI}
LangevinLean Project. \emph{Unadjusted Langevin under LSI}. LL-02, 2026. Draft report, content revision 963f75cb40db. \url{https://odakyildiz.github.io/LangevinLean/docs/pdf/UnadjustedLangevin/LSI/FullProof.pdf}.
\bibitem[IPLA-SC]{LangevinLean2026IPLA}
LangevinLean Project. \emph{Interacting Particle Langevin under Strong Convexity}. LL-03, 2026. Draft report, content revision ec0201df2207. \url{https://odakyildiz.github.io/LangevinLean/docs/pdf/IPLA/FullProof.pdf}.
\end{thebibliography}

\end{document}
